% chapters/13/section-05.tex
\section{Masalah Fisika dan Perbaikannya}
\label{sec:masalah-fisika}

Inilah bagian inti bab ini, yaitu menganalisis dan memperbaiki masalah fisika,
sesuai indikator \ind{c}. Empat masalah berikut paling sering terjadi, dan
keempatnya memiliki gejala yang khas.

\subsection{Empat Masalah yang Paling Sering}

\begin{table}[htbp]
\centering
\caption{Empat masalah fisika, penyebabnya, dan perbaikannya}
\label{tab:masalah-fisika}
{\setlength{\tabcolsep}{4pt}\footnotesize%
\begin{tabularx}{\textwidth}{@{}L{2.8cm} L{4.2cm} Y@{}}
\toprule
\headrow \textbf{Gejala} & \textbf{Penyebab} & \textbf{Perbaikan} \\
\midrule
Benda bergetar ketika bergerak atau ketika kamera mengikutinya &
Benda digerakkan melalui \code{transform.position} pada \code{Update},
sehingga letaknya berubah di antara dua langkah fisika. &
Gerakkan melalui \istilah{Rigidbody} pada \code{FixedUpdate}, dan aktifkan
\istilah{Interpolate} pada benda yang diikuti kamera. \\
\addlinespace
Benda menembus dinding atau lantai, terutama ketika bergerak cepat &
Pemeriksaan tabrakan bawaan dilakukan pada saat-saat tertentu saja, sehingga
benda yang cepat dapat melompati benda lain di antaranya. &
Ubah \istilah{Collision Detection} menjadi \istilah{Continuous}, atau
kurangi kecepatannya. Gejala ini sama dengan \istilah{tunneling} pada
Bagian~2.4. \\
\addlinespace
Benda jatuh tanpa henti &
Benda tidak memiliki \istilah{Collider}, atau lantainya tidak memiliki
\istilah{Collider}, atau keduanya berada pada \istilah{layer} yang tidak
saling diperiksa. &
Periksa ketiga hal itu berturut-turut, dengan menampilkan bentuk batas benda
melalui panel \istilah{Physics Debug}. \\
\addlinespace
Benda tersangkut pada dinding, atau memantul aneh pada sudut &
Benda memiliki dua \istilah{Collider} yang saling bertumpang tindih, atau
memakai \istilah{Collider} dari gambar yang bentuknya rumit, atau bergesekan
dengan dinding tanpa bahan gesekan. &
Pakai satu \istilah{Collider} sederhana, dan berikan bahan fisika dengan
gesekan nol pada benda yang bergesekan dengan dinding. \\
\bottomrule
\end{tabularx}}
\end{table}

\subsection{Menampilkan Bentuk Batas Benda}

Cara paling cepat untuk memeriksa masalah fisika adalah dengan
\textbf{menampilkan bentuk batasnya}. Apabila bentuk batas itu tidak tampak
pada tempat yang seharusnya, penyebabnya hampir pasti berkaitan dengan letak
atau ukuran \istilah{Collider}, bukan dengan kode Anda.

Panel \istilah{Physics Debug} pada tangkapan layar
Gambar~\ref{fig:editor-fisika} memperlihatkan hal itu: kotak kuning pada gambar
tersebut adalah bentuk batas sebuah benda, dan terlihat bahwa bentuknya jauh
lebih sederhana daripada gambarnya. Panel itu dibuka melalui menu
\istilah{Window}, lalu \istilah{Analysis}, lalu \istilah{Physics Debug}.

\begin{tipsbox}
Kebiasaan yang menghemat waktu: apabila perilaku fisika terasa janggal,
tampilkan lebih dahulu bentuk batas benda, dan periksa apakah letak serta
ukurannya sudah sesuai dugaan. Sebagian besar masalah fisika pada tahap
pembelajaran berasal dari \istilah{Collider}, bukan dari kode.
\end{tipsbox}

\subsection{Urutan Penelusuran yang Disarankan}

Empat langkah berikut disarankan untuk menelusuri masalah fisika, dan
urutannya sengaja disusun dari yang paling murah.

\begin{langkah}
  \item \textbf{Periksa komponennya pada \istilah{Inspector}.} Apakah benda
        yang bergerak memiliki \istilah{Rigidbody} dan \istilah{Collider}?
        Apakah benda diam cukup memiliki \istilah{Collider}? Apakah
        \istilah{Collider} pada benda yang bergerak ditandai sebagai
        \istilah{trigger} tanpa sengaja?
  \item \textbf{Tampilkan bentuk batas benda} melalui \istilah{Physics Debug},
        lalu periksa letak dan ukurannya.
  \item \textbf{Periksa pengaturan fisika} yang berkaitan dengan gejala:
        \istilah{Gravity Scale} untuk masalah berat, \istilah{Collision
        Detection} untuk masalah penembusan, dan \istilah{Interpolate} untuk
        masalah getaran.
  \item \textbf{Baru periksa kode Anda}, dan periksa lebih dahulu apakah
        \istilah{Rigidbody} digerakkan pada fungsi yang tepat.
\end{langkah}

\begin{peringatanbox}
Perhatikan urutan itu. Pada blok Python, masalah selalu berada pada kode Anda,
sebab seluruh perhitungannya memang Anda tulis sendiri. Pada Unity, sebagian
besar masalah justru berada pada pengaturan komponen. Menelusuri kode terlebih
dahulu merupakan kebiasaan dari blok sebelumnya yang perlu disesuaikan.
\end{peringatanbox}
